% 原课程信息. 在 \frontmatter 内输入, 不增加数学章号或数学环境编号.
\par\addvspace{14pt}\Needspace{5\baselineskip}
\phantomsection\addcontentsline{toc}{chapter}{原课程信息与学习安排}
\markboth{原课程信息与学习安排}{}
\noindent{\color{structurecolor}\bfseries\fontsize{14.4}{20}\selectfont 原课程信息与学习安排}\par\nobreak\vspace{6pt}

以下分别保留本书所用八门 MIT OpenCourseWare 课程的原始身份和教学安排. 18.S190 是度量空间部分的主源; 18.901 支持一般拓扑主线; 18.900 支持组合曲面部分. 其余五门课程用于补充路线或选读材料. 原课先修与评分属于各自学期, 不构成本书的统一先修要求、考纲或评分制度. 原课讲次与本书章次也不相同.

课程信息据官方主页、Syllabus、Calendar 及相关材料页核对. 有些页面只列讲次和作业节点, 没有实际日期或星期; 以下按原有粒度转述. 对未在所用公开页面列出的信息明确说明, 不依据每周课时数倒推课表. 本书的证明补全、中文解答与编者练习另见正文和来源说明.

\par\addvspace{9pt}\Needspace{4\baselineskip}
\phantomsection\addcontentsline{toc}{section}{18.S190: 度量空间入门}
\noindent{\bfseries\fontsize{12}{17}\selectfont 18.S190: 度量空间入门}\par\nobreak\vspace{4pt}

\textbf{作者、学期与角色.} Introduction to Metric Spaces, January IAP 2023, 授课者 Paige Bright, 本科课程. 本书第1--8章以六讲公开讲义为主要起点, 按概念依赖重编并补全论证. 原课介绍度量、开闭集、连续映射、函数空间、完备性与紧致性, 是短期专题课, 并非一般拓扑的全学期课程.

\textbf{上课安排与先修.} 每周2次讲课, 每次1.5小时. 官方先修为18.100A Real Analysis 或18.100P Real Analysis. 这说明原课面向已接触分析和证明写作的学生; 本书另在前言说明自身读者要求.

\textbf{作业、考试与评分.} 原课采用 Pass/No Record, 依据作业与参与情况评定; 共3份 problem sets, 无考试. 官方没有列出作业与参与的百分比, 不应另配权重.

\textbf{官方 Calendar.} 第1讲为动机、直观与例子; 第2讲为一般理论, 交第1份作业; 第3讲为欧氏空间中的紧集, 交第2份作业; 第4讲为紧度量空间; 第5讲为完备度量空间, 交第3份作业; 第6讲为后续方向. Calendar 位于 Syllabus 内, 只列讲次及上述提交节点, 未列实际日期.

\textbf{依据.} \href{https://ocw.mit.edu/courses/18-s190-introduction-to-metric-spaces-january-iap-2023/}{官方课程主页}; \href{https://ocw.mit.edu/courses/18-s190-introduction-to-metric-spaces-january-iap-2023/pages/syllabus/}{Syllabus 与 Calendar}.

\par\addvspace{9pt}\Needspace{4\baselineskip}
\phantomsection\addcontentsline{toc}{section}{18.901: 一般拓扑主线}
\noindent{\bfseries\fontsize{12}{17}\selectfont 18.901: 一般拓扑主线}\par\nobreak\vspace{4pt}

\textbf{作者、学期与角色.} Introduction to Topology, Fall 2004, 授课者 James Munkres, 本科课程. 原课以逻辑与集合论、一般拓扑和专业数学写作为目标. 教材为 Munkres 的 \emph{Topology}, 第2版; OCW 提供补充讲义和按教材定位的练习, 不应把这些文件称作全课完整公开讲义. 本书第9--10章采用其一般拓扑安排及补充 Notes D、G, 完整教学论证由编者组织和补充.

\textbf{上课安排与先修.} 每周2次讲课, 每次1.5小时. 原课要求第一门分析课的水平达到 Rudin 的 \emph{Principles of Mathematical Analysis}, 在 MIT 对应18.100B; 既要求分析知识, 也要求构造和书写证明的经验. 每周在课堂讨论学生希望提问的习题.

\textbf{周练与正式作业.} 每次讲课配套的教材习题用于训练, 不提交评分; 原课鼓励把全部解答写成笔记, 考试几乎全部以这些习题为基础. 另有4份正式 problem sets, 须独立完成、提交并按数学写作要求评分. PS0 是集合论语言的诊断作业, 第二次课在课堂评阅, 成绩仅用于学习建议; PS5 是学期末的可选作业. 不应把编号0--5的六组材料都计作正式评分作业.

\textbf{评分.} 正式作业共400分, 期中考试100分, 期末考试200分, 总计700分. Syllabus 未把它们写成百分制比例; 本书保留原分值.

\textbf{官方 Calendar: 第1--5周.} 第1周为逻辑与基础(第1次课), 列 PS0 提交节点; 第2周为关系、基数与选择公理(第2--3次课); 第3周为拓扑与闭集(第4--5次课); 第4周为连续函数与任意积(第6--7次课), 在该周第二次课交 PS1; 第5周为度量拓扑(第8--9次课).

\textbf{官方 Calendar: 第6--10周.} 第6周为商拓扑(第10次课); 第7周为连通空间与紧空间(第11--12次课); 第8周继续紧致性(第13--14次课), 在第二次课交 PS2; 第9周先讲良序集与极大原理(第15次课), 第16次课为期中考试; 第10周为可数性与分离公理(第17--18次课).

\textbf{官方 Calendar: 第11周以后.} 第11周为 Urysohn 引理与度量化(第19--20次课), 在第二次课交 PS3; 第12周为 Tietze 定理(第21次课); 第13周为 Tychonoff 定理与 Stone--\v{C}ech 紧化(第22--23次课); 第14周为 Baire 空间与维数理论(第24--25次课), 在第二次课交 PS4; 第15周为嵌入欧氏空间(第26次课). 最后一行列期末考试和可选 PS5 的提交, 未给该行周号或实际日期.

\textbf{依据.} \href{https://ocw.mit.edu/courses/18-901-introduction-to-topology-fall-2004/}{官方课程主页}; \href{https://ocw.mit.edu/courses/18-901-introduction-to-topology-fall-2004/pages/syllabus/}{Syllabus}; \href{https://ocw.mit.edu/courses/18-901-introduction-to-topology-fall-2004/pages/calendar/}{Calendar}; \href{https://ocw.mit.edu/courses/18-901-introduction-to-topology-fall-2004/pages/assignments/}{Assignments}.

\par\addvspace{9pt}\Needspace{4\baselineskip}
\phantomsection\addcontentsline{toc}{section}{18.900: 平面几何与拓扑}
\noindent{\bfseries\fontsize{12}{17}\selectfont 18.900: 平面几何与拓扑}\par\nobreak\vspace{4pt}

\textbf{作者、学期与角色.} Geometry and Topology in the Plane, Spring 2023, 授课者 Paul Seidel, 本科课程. 原课选取易于可视化的几何与拓扑专题, 不要求已有微分几何或拓扑学训练. 本书使用第4讲的绕数相关材料及第29--31、33、40讲的组合曲面材料; 这只是原课的一部分, 不包含其台球、代数曲线或双曲几何等完整专题.

\textbf{上课安排与先修.} 每周3次讲课, 每次50分钟. 官方先修为18.03 Differential Equations 或18.06 Linear Algebra, 并预期熟悉18.02 Multivariable Calculus. 原课解释线性代数在第III章短暂出现, 在第VII章作用更大; 微分方程用于第IX章, 复数用于第14讲及第VIII章.

\textbf{作业、考试与评分.} 理解性问题(comprehension questions, 简称 CQ)占30\%, 较有挑战的正式 problem sets 占30\%, 考试占40\%. CQ 在 OCW 公开, 正式 problem sets 和考试不向 OCW 用户提供. 因而不能从网页的 Problem Sets 资源标签推断正式评分题或试卷已公开, 本书所写解答也不冒称这些未公开材料的官方答案.

\textbf{公开讲义顺序与实际授课.} 公开目录有9章、40讲: 第1--5讲为多边形, 第6--8讲为台球, 第9--11讲为主频, 第12--15讲为光滑闭路, 第16--18讲为浸入闭路, 第19--27讲为代数曲线, 第28--33讲为二维复形, 第34--36讲为双曲几何, 第37--40讲为弯曲几何. 这是可用材料的顺序; Syllabus 明确说明2023年春实际略过第20、24--25、38和40讲. 本书第14章采用第40讲 \emph{Geometry of Combinatorial Surfaces} 作为 OCW 公开选读, 不把它标为该学期实际讲授内容.

\textbf{官方 Calendar.} 独立的 Calendar 及带实际日期的授课、作业或考试时间表未在所用公开页面列出. 上述40讲目录与实际略讲说明可以确定材料范围, 不能据此合成逐周课表.

\textbf{依据.} \href{https://ocw.mit.edu/courses/18-900-geometry-and-topology-in-the-plane-spring-2023/}{官方课程主页与讲义目录}; \href{https://ocw.mit.edu/courses/18-900-geometry-and-topology-in-the-plane-spring-2023/pages/syllabus/}{Syllabus}.

\par\addvspace{9pt}\Needspace{4\baselineskip}
\phantomsection\addcontentsline{toc}{section}{18.904: 拓扑研讨课}
\noindent{\bfseries\fontsize{12}{17}\selectfont 18.904: 拓扑研讨课}\par\nobreak\vspace{4pt}

\textbf{作者、学期与角色.} Seminar in Topology, Spring 2011, 授课者 Andrew Snowden, 本科课程. 学生每次课分两人讲授, 每人约25分钟; 目标同时包括基本群、同调、上同调和数学报告能力. 本书第11--13章参照公开 Lecture Summaries 的基本群与覆盖路线, 摘要不等同于完整原讲义.

\textbf{上课安排、先修与评分.} 每周3次, 每次1小时; 先修18.901 Introduction to Topology. 讲授与参与占60\%, 期末论文占30\%, 作业占10\%; 无考试. Syllabus 预计约4份正式作业, 讲者提供的练习集可选且不提交. 须出勤, 每缺3次课降低一个字母等级; 期末论文至少10页, 用 LaTeX 编写.

\textbf{官方 Calendar.} 第1--14次课为基本群与覆盖, PS1 于第9次课提交; 第15--25次课为同调, 第15次课定论文题目, PS2 于第20次课提交; 第26--34次课为上同调与 Poincar\'e 对偶, 第27次课交论文初稿、第32次课交终稿; 第35--39次课为期末论文报告, PS3 于第38次课提交. Calendar 只明确列出这3份作业节点, 不应为预计的第4份补造时间.

\textbf{依据.} \href{https://ocw.mit.edu/courses/18-904-seminar-in-topology-spring-2011/}{官方课程主页}; \href{https://ocw.mit.edu/courses/18-904-seminar-in-topology-spring-2011/pages/syllabus/}{Syllabus 与 Calendar}.

\par\addvspace{9pt}\Needspace{4\baselineskip}
\phantomsection\addcontentsline{toc}{section}{18.905: 代数拓扑选读}
\noindent{\bfseries\fontsize{12}{17}\selectfont 18.905: 代数拓扑选读}\par\nobreak\vspace{4pt}

\textbf{作者、学期与角色.} Algebraic Topology I, Fall 2016, 授课者 Haynes Miller, 研究生课程. 官方说明公开讲义依据 Sanath Devalapurkar 的课堂 LaTeX 记录, 插图由 Xianglong Ni 绘制. 本书选用第5讲同伦、第14讲 CW 复形及第18讲 Euler 示性数材料. 原课从奇异同调起步, 不能当作从零讲授基本群与覆盖的课程.

\textbf{上课安排、先修与评分.} 每周2次, 每次1小时; 先修为18.701 Algebra I 或18.703 Modern Algebra, 加18.901 Introduction to Topology. 须熟悉拓扑空间、覆盖和基本群, 以及主理想整环上有限生成模的结构. 共6份作业, 占75\%; 期末考试周有40分钟口试, 占25\%.

\textbf{官方 Calendar.} 第1--13讲为奇异同调, 第14--25讲为计算方法, 第26--38讲为上同调与对偶. PS1--6 分别于第7、12、16、22、28、35讲提交; 口试安排在期末考试周. 页面按讲次列出节点, 未列实际日期.

\textbf{依据.} \href{https://ocw.mit.edu/courses/18-905-algebraic-topology-i-fall-2016/}{官方课程主页}; \href{https://ocw.mit.edu/courses/18-905-algebraic-topology-i-fall-2016/pages/syllabus/}{Syllabus}; \href{https://ocw.mit.edu/courses/18-905-algebraic-topology-i-fall-2016/pages/calendar/}{Calendar}; \href{https://ocw.mit.edu/courses/18-905-algebraic-topology-i-fall-2016/pages/lecture-notes/}{讲义署名与目录}.

\par\addvspace{9pt}\Needspace{4\baselineskip}
\phantomsection\addcontentsline{toc}{section}{18.950: 微分几何选读}
\noindent{\bfseries\fontsize{12}{17}\selectfont 18.950: 微分几何选读}\par\nobreak\vspace{4pt}

\textbf{作者、学期与角色.} Differential Geometry, Fall 2008, 授课者 Paul Seidel, 本科课程. 原课围绕曲率, 讨论平面曲线、超曲面的局部及整体几何、长度与距离. 本书第15章只选取公开第4章的长度与距离材料, 不据此宣称完整覆盖18.950.

\textbf{上课安排与先修.} 每周上课次数和每次时长未在所用公开页面列出. 官方先修为18.100 Analysis I, 加18.06或18.700 Linear Algebra 或18.701 Algebra I 中的一门; 课程说明同时要求良好的多元微积分和线性代数基础, 并适应定义、定理与证明的叙述.

\textbf{评分与公开作业.} 期中考试占30\%, 作业占30\%, 期末考试占40\%. Assignments 页公开 Homework 1--10. 十份作业的中文题面、独立解答及覆盖记录见本项目独立分册\href{https://github.com/xhc144/mit-ocw-notes/blob/80b9182650861c1a4830fb812032abbc93a429e7/subjects/differential-geometry/differential-geometry.pdf}{《曲线与曲面的微分几何》PDF}, 可下载其\href{https://github.com/xhc144/mit-ocw-notes/blob/80b9182650861c1a4830fb812032abbc93a429e7/subjects/differential-geometry/differential-geometry-source.zip}{源码 ZIP}. 该分册说明32个顶层题中30题的全部要求已有独立解答, PS3.3末尾引用的 Proposition 6.3 与 PS9.4引用的 Lemma 28.3 尚有定位缺口. 本册不重复这些作业解答.

\textbf{官方 Calendar.} 独立 Calendar 及作业、考试实际日期未在所用公开页面列出. Lecture Notes 目录仅给讲次分段: 第1--10讲为平面曲线的局部与整体几何, 第11--23讲为超曲面局部几何, 第24--35讲为超曲面整体几何, 第36--41讲为长度与距离. 这不是带日期的周课表.

\textbf{依据.} \href{https://ocw.mit.edu/courses/18-950-differential-geometry-fall-2008/}{官方课程主页}; \href{https://ocw.mit.edu/courses/18-950-differential-geometry-fall-2008/pages/syllabus/}{Syllabus}; \href{https://ocw.mit.edu/courses/18-950-differential-geometry-fall-2008/pages/lecture-notes/}{讲义目录}; \href{https://ocw.mit.edu/courses/18-950-differential-geometry-fall-2008/pages/assignments/}{10份作业}; \href{https://github.com/xhc144/mit-ocw-notes/blob/80b9182650861c1a4830fb812032abbc93a429e7/subjects/differential-geometry/README.md}{微分几何分册说明}.

\par\addvspace{9pt}\Needspace{4\baselineskip}
\phantomsection\addcontentsline{toc}{section}{18.965: 流形几何选读}
\noindent{\bfseries\fontsize{12}{17}\selectfont 18.965: 流形几何选读}\par\nobreak\vspace{4pt}

\textbf{作者、学期与角色.} Geometry of Manifolds, Fall 2004, 授课者 Tomasz Mrowka, 研究生课程. 本书第15章选用第1、2、4讲的基础坐标材料, 不包含原课的全部微分拓扑理论.

\textbf{上课安排、先修与评分.} 每周2次, 每次1.5小时; 先修18.101 Analysis II 与18.905 Algebraic Topology. 成绩100\%依据作业. 所用 Syllabus 未另列考试权重.

\textbf{官方 Calendar.} 共列35讲, 从流形、光滑映射、导数和逆函数及隐函数定理出发, 进入向量丛、嵌入、Sard 定理、纤维丛、Fredholm 算子与横截性. 第23--28讲为 Morse 理论; 第29--30讲涉及 Lie 导数和 Frobenius 可积性; 第31--34讲为微分形式及 de Rham 定理相关内容; 第35讲为 Smale 浸入定理. Calendar 未列作业提交或考试日期.

\textbf{依据.} \href{https://ocw.mit.edu/courses/18-965-geometry-of-manifolds-fall-2004/}{官方课程主页}; \href{https://ocw.mit.edu/courses/18-965-geometry-of-manifolds-fall-2004/pages/syllabus/}{Syllabus}; \href{https://ocw.mit.edu/courses/18-965-geometry-of-manifolds-fall-2004/pages/calendar/}{Calendar}.

\par\addvspace{9pt}\Needspace{4\baselineskip}
\phantomsection\addcontentsline{toc}{section}{18.102: 泛函分析补充}
\noindent{\bfseries\fontsize{12}{17}\selectfont 18.102: 泛函分析补充}\par\nobreak\vspace{4pt}

\textbf{作者、学期与角色.} Introduction to Functional Analysis, Spring 2021, 授课者 Casey Rodriguez, 本科课程. 官方 Lecture Notes and Readings 页区分 Richard Melrose 的2020年讲义与 Andrew Lin 在 Rodriguez 2021年课堂所作的记录; 本书第8章使用后者第3讲中的 Baire 定理与一致有界原理. 该补充支持分析应用, 不代表本册包含原课的测度论、Hilbert 空间或谱定理全程.

\textbf{上课安排与先修.} 每周2次, 每次1.5小时. 先修为18.06、18.700或18.701之一, 加18.100A、18.100B、18.100P或18.100Q之一; 或获教师许可.

\textbf{作业、考试与评分.} 每周作业占50\%, 去掉最低一次作业成绩; 期中考试占25\%, 为第7周的24小时带回考试; 期末 assignment 占25\%, 为课程结束时的48小时作业. 原页将后者称为 Final Assignment, 本书不改称正式期末考试.

\textbf{官方 Calendar.} 共列23讲: 第1--5讲建立 Banach 空间、算子、Baire 及 Hahn--Banach 等基础; 第6--13讲转入测度、Lebesgue 积分与 $L^p$ 空间; 第14--22讲为 Hilbert 空间、紧算子与谱理论; 第23讲为区间上的 Dirichlet 问题. Assignment 1--10 分别于第2、4、6、8、10、14、16、18、20、22讲提交. 期中考试列在第11讲之后、第12讲之前, 最后一行列 Final Assignment 提交; 页面未列实际日期.

\textbf{依据.} \href{https://ocw.mit.edu/courses/18-102-introduction-to-functional-analysis-spring-2021/}{官方课程主页}; \href{https://ocw.mit.edu/courses/18-102-introduction-to-functional-analysis-spring-2021/pages/syllabus/}{Syllabus}; \href{https://ocw.mit.edu/courses/18-102-introduction-to-functional-analysis-spring-2021/pages/calendar/}{Calendar}; \href{https://ocw.mit.edu/courses/18-102-introduction-to-functional-analysis-spring-2021/pages/lecture-notes-and-readings/}{讲义署名与阅读安排}.
